% Created 2026-06-27 Sat 15:17
% Intended LaTeX compiler: pdflatex
\documentclass[11pt,letterpaper]{article}
\usepackage[utf8]{inputenc}
\usepackage[T1]{fontenc}
\usepackage{graphicx}
\usepackage{longtable}
\usepackage{wrapfig}
\usepackage{rotating}
\usepackage[normalem]{ulem}
\usepackage{amsmath}
\usepackage{amssymb}
\usepackage{capt-of}
\usepackage{hyperref}
\usepackage{geometry}
\geometry{left=1.2in, right=1.2in, top=1.0in, bottom=1.0in}
\usepackage{fancyhdr}
\usepackage{parskip}
\usepackage{microtype}
\usepackage{booktabs}
\usepackage{enumitem}
\usepackage[utf8]{inputenc}
\usepackage[T1]{fontenc}
\usepackage{textgreek}
\PassOptionsToPackage{hidelinks}{hyperref}
\setlength{\parindent}{0pt}
\setlength{\parskip}{6pt}
\fancypagestyle{plain}{\fancyhf{}\fancyfoot[C]{\thepage}\renewcommand{\headrulewidth}{0pt}}
\fancypagestyle{body}{%
\fancyhf{}%
\fancyhead[C]{\small The Medium --- Action Plan $\cdot$ 2026}%
\fancyfoot[C]{\thepage}%
\renewcommand{\headrulewidth}{0.4pt}}
\pagestyle{body}
\author{Bradley K. DePuy}
\date{2026-06-27}
\title{The Medium --- Action Plan: Open Problems and Priorities}
\hypersetup{
 pdfauthor={Bradley K. DePuy},
 pdftitle={The Medium --- Action Plan: Open Problems and Priorities},
 pdfkeywords={},
 pdfsubject={},
 pdfcreator={Emacs 29.3 (Org mode 9.6.15)}, 
 pdflang={English}}
\begin{document}

\maketitle
\tableofcontents


\section{Overview}
\label{sec:org714e66d}

This document replaces the original Smolin/Rovelli/Hossenfelder critical review
action plan. Many issues from that review have been addressed in the current
document revisions. This plan tracks what remains open, what has been resolved,
and the priority order for future work.

\subsection{What Has Been Resolved}
\label{sec:org0d64acc}

The following issues from the original review have been addressed in the
current document set (docs 0--5):

\begin{itemize}
\item \textbf{SU(3) claim corrected} --- body and abstract now correctly state U(1) ×U(2)
as the established group; SU(3) moved to open derivation status (Doc 3)
\item \textbf{LQG minimum area match removed} --- 4π discrepancy between The Medium and
LQG formulas now explicitly noted (Doc 4)
\item \textbf{Diffeomorphism invariance claim hedged} --- "inherits" replaced with
structural motivation shared; constraint algebra closing identified as open (Doc 4)
\item \textbf{Proton mass labeling fixed} --- "Exact" changed to "Exact (by construction)"
in all tables (Docs 1, 2, 4, 5)
\item \textbf{G and m\_p variation predictions flagged} --- moved from predbox to openbox;
explicitly noted as appearing observationally ruled out with specific numbers (Doc 5)
\item \textbf{H\(_0\) predbox softened} --- "Prediction" changed to "Interpretation"; selected
value vs. derived value distinction made clear (Doc 5)
\item \textbf{Ω\_Λ = 2/3 relabeled} --- "Prediction" changed to "Consequence of G formula";
not independently testable (Doc 5)
\item \textbf{Λ hierarchy claim made conditional} --- floor derivation identified as
required; "dissolving" language removed (Docs 3, 5)
\item \textbf{ADM summary corrected} --- "recovered" replaced with schematic structure
present; missing components listed (Doc 4)
\item \textbf{Pion added to hadron table} --- explicitly noted as not explained by current
formula (Doc 2)
\item \textbf{"Complete chain" removed} --- changed to "connected chain\ldots{} with open
calculations explicitly identified" (Doc 2)
\item \textbf{Spin identification qualified} --- SU(2) representation theory noted as not
yet derived (Doc 4)
\item \textbf{Color charge claim qualified} --- "are color charge" changed to "structurally
analogous to color charge labels" (Doc 2)
\end{itemize}

\section{{\bfseries\sffamily FATAL} Issues}
\label{sec:org5f0cef4}

\subsection{{\bfseries\sffamily FATAL} [F1] Foundation connection --- substantially addressed, four gaps remain}
\label{sec:org74fc5dd}
Doc 1 has been substantially revised. The following have been resolved:

\begin{itemize}
\item[{$\boxtimes$}] \(M\) defined: \(M\) is the substrate itself. Not a geometric space.
\item[{$\boxtimes$}] \(\iota\) is a distinction within \(M\), not an event (no sequence implied).
\item[{$\boxtimes$}] \(\phi: M \to \{+1,-1\}\) replaced by \(c: D \times D \to \{\mathrm{same}, \mathrm{different}\}\). Character is relational, not intrinsic.
\item[{$\boxtimes$}] Relations between co-existing distinctions are unavoidable (self-grounding argument: follows from singularity of \(M\), same logic as distinction itself requiring no cause).
\item[{$\boxtimes$}] Adjacency constraint reframed as informational, not topological.
\item[{$\boxtimes$}] Closure condition identified as the localization mechanism: closed structures are self-contained; open chains are unstable and must close.
\item[{$\boxtimes$}] Degree-2 argument stated and its implication documented (see below).
\end{itemize}

Four gaps remain open:

\subsubsection{Gap 1 --- Degree-2 not yet proven}
\label{sec:org7af9d59}

The argument that each \(\iota\) distinction has relational degree exactly 2 (from the two sides of its split) is plausible but not rigorous. Nothing in the current formulation rules out higher-degree nodes. If degree is not fixed at 2, the triangle is not uniquely minimal and the localization argument weakens.

\textbf{Testable implication:} If degree 2 is correct, \(K_3\) is the only primitive closed structure. All higher configurations (\(K_4\), \(K_5\)) must be composite. This is consistent with QCD but must be verified against the full particle spectrum.

\textbf{Action:} Derive degree-2 from the binary split rigorously, or identify the additional structure that allows higher degree.

\subsubsection{Gap 2 --- Mathematical construction of N from c not specified}
\label{sec:org72179f7}

\(N_{ij}\) is described as "the coherence of comparison across all paths through \(G\)." The precise construction --- how the binary comparison \(c(i,j) \in \{\mathrm{same}, \mathrm{different}\}\) gives rise to the continuous real field \(N_{ij} \in \mathbb{R}\) --- is not written down. The candidate identification \(N_{ij} \sim (L^{-1})_{ij}\) is plausible but circular (L depends on G, G depends on N).

\textbf{Action:} Write an explicit non-circular construction of \(N\) from \(c\).

\subsubsection{Gap 3 --- Complex N has no source}
\label{sec:orgfd83474}

The full Lagrangian and any route to SU(3) symmetry requires \(N \in \mathbb{C}\). Neither the binary comparison \(c\) nor the real-valued coherence argument gives a natural origin for complex phases. Promoting \(c\) to a U(1)-valued comparison (phases rather than binary values) is the most natural candidate but has not been developed.

\textbf{Action:} Identify the mechanism that produces complex phase in \(N\).

\subsubsection{Gap 4 --- Frustrated triangle ground state not worked out}
\label{sec:org69073a7}

The triangle \(K_3\) necessarily has at least one same-comparing edge (it cannot be 2-colored without adjacent same-color nodes). The adjacency constraint says same-comparing negotiations are informationally empty. What is \(N\) on the frustrated edge? If \(N = 0\), the triangle loses an edge and the Laplacian changes. If \(N \neq 0\), the constraint is violated. The degenerate ground state of the frustrated triangle --- from which spin is supposed to emerge --- has not been derived.

\textbf{Action:} Work out the ground state of the frustrated triangle under the informational adjacency constraint. Determine whether frustration degeneracy corresponds to the double degeneracy of \(\lambda = 3\).

\subsection{{\bfseries\sffamily RESOLVED} [F2] "From First Principles" on Doc 2 title page}
\label{sec:org8b3e677}
\begin{itemize}
\item[{$\boxtimes$}] Subtitle changed to "Structure from Substrate Primitives; Scale from Observation."
\end{itemize}

The distinction: structural results (triangle topology, \(\lambda = 3\), three-fold
mass ratio, quantization from closure, confinement from self-containment) are
genuinely derived from the substrate primitives. The numerical scale (\(\varepsilon\),
\(H_0\), \(\gamma\)) requires empirical input. The new subtitle reflects this
accurately.

\section{{\bfseries\sffamily SEVERE} Issues}
\label{sec:orge4253ed}

\subsection{{\bfseries\sffamily SEVERE} [S1] G and m\_p variation rates appear observationally ruled out}
\label{sec:org008b80a}
The cosmological evolution equations give:
\begin{itemize}
\item \(\dot{G}/G \approx -1.5 \times 10^{-10}\)/yr
\item \(\dot{m}_p/m_p \approx 7 \times 10^{-11}\)/yr
\end{itemize}

Current observational bounds are approximately 200× tighter for \(G\) and approximately 7000× tighter
for \(m_p\). These are now flagged in Doc 5 as "may be observationally ruled out"
in an openbox. This is correct but incomplete --- the underlying equations need
to be examined.

\subsubsection{Action steps}
\label{sec:org9be7d32}

\begin{enumerate}
\item{$\square$} Determine whether a physical mechanism exists that suppresses local
variation while allowing secular variation over Hubble timescales. This is
the only route to saving these predictions. Candidates: the floor condition
at \(\rho_\Lambda\), or a relational definition of \(G\) and \(m_p\) that
makes them observer-dependent rather than absolute.
\item{$\square$} If no such mechanism can be identified, remove the evolution equations
for \(G\) and \(m_p\) from the description entirely, or relocate them to
a section labeled "candidate equations pending constraint analysis."
\item{$\square$} \(\Lambda\) variation (\(d\Lambda/dt = -2\Lambda H_0\)) has no current
constraint but follows from the same equations. Its status is linked to
items 1 and 2.
\end{enumerate}

\subsection{{\bfseries\sffamily SEVERE} [S2] Fermionic statistics not produced by the bosonic Lagrangian}
\label{sec:org76ca0ad}
The negotiation Lagrangian is bosonic. Canonical quantization gives
\([\hat{a}, \hat{a}^\dagger] = 1\) (Bose-Einstein statistics). The proton is
a spin-1/2 fermion. This is identified as a primary open problem in Docs 3 and 4
but no progress has been made on either candidate resolution.

\subsubsection{Action steps}
\label{sec:orgca05ab8}

\begin{enumerate}
\item{$\square$} Investigate composite antisymmetrization first: show whether three
bosonic quanta in the antisymmetric combination on the triangle nodes acquire
fermionic exchange behavior under interchange of two protons. Write out the
explicit state and check the exchange statistics.
\item{$\square$} If antisymmetrization fails: investigate Grassmann-valued \(N\) and
rederive the Lagrangian with anticommuting fields.
\item{$\square$} Document whichever path is pursued, even if it fails, so a mathematician
reading this knows what has been tried.
\end{enumerate}

\subsection{{\bfseries\sffamily SEVERE} [S3] The energy formula E(n,k) = ε·k·n is not fully derived}
\label{sec:org5eca076}
The formula is presented as following from the Hamiltonian eigenvalue equation,
and for complete graphs \(K_n\) with mode \(k=1\) (ground non-trivial mode),
\(E = \varepsilon \cdot n\) follows directly from \(\lambda = n\). But the \(k\)
factor for higher modes is introduced without derivation, and the table in Doc 4
uses \(E_n = \varepsilon \cdot n\) without specifying \(k\). For the rho meson
(n=2) the predicted value of 625.6 MeV already differs from observed (775 MeV)
by 19\%. The formula's domain of validity is unclear.

\subsubsection{Action steps}
\label{sec:org7934fd4}

\begin{enumerate}
\item{$\square$} Derive explicitly what \(k\) represents for higher negotiation modes on
complete graphs and write out \(E(n,k)\) with \(k\) specified for each entry
in the hadron table.
\item{$\square$} Address the rho meson 19\% discrepancy: is the formula valid for n=2
or does the two-node topology require a different treatment?
\end{enumerate}

\section{{\bfseries\sffamily SERIOUS} Issues}
\label{sec:orgc9caffd}

\subsection{{\bfseries\sffamily SERIOUS} [SR1] The symmetry group of the triangle Lagrangian not established}
\label{sec:org8468c82}
Doc 3 now correctly states that the group is U(1) ×U(2), not SU(3), and
identifies SU(3) as an open derivation. But the path from U(1) ×U(2) to SU(3)
--- if it exists --- has not been investigated. This blocks the color charge
identification and the strong force correspondence.

\subsubsection{Action steps}
\label{sec:orgf6b7ff9}

\begin{enumerate}
\item{$\square$} Compute the symmetry group of \(N^\dagger L N\) for \(K_3\) with complex
\(N\) explicitly. Confirm U(1) ×U(2).
\item{$\square$} Identify what additional invariant --- a cubic \(\varepsilon_{ijk}\) term,
a different action on edge states --- could promote the symmetry to SU(3).
\item{$\square$} If SU(3) cannot be reached: accept U(1) ×U(2) as the established group
and demote all color charge correspondence to "structural analogy at the level
of node counting."
\end{enumerate}

\subsection{{\bfseries\sffamily SERIOUS} [SR2] The pion is absent from the hadron spectrum}
\label{sec:orgb5c215c}
The pion (135 MeV) is the lightest hadron and the most important test case.
It is now explicitly listed in the Doc 2 hadron table as "not explained by this
formula." A physicist reviewing this will treat the pion's absence as a primary
failure of the spectrum description, not a minor gap. It requires a qualitatively
different physical mechanism (chiral symmetry breaking) that is not present in
the current description.

\subsubsection{Action steps}
\label{sec:org02f87bc}

\begin{enumerate}
\item{$\square$} Investigate whether the pion can be accommodated as a two-node graph with
a non-complete topology (i.e., a single edge rather than the complete graph),
or as a Goldstone mode of a broken symmetry in the negotiation field.
\item{$\square$} If neither works: explicitly state in Doc 2 that the light meson sector
requires a mechanism outside the current description, and identify what that
mechanism would need to be.
\end{enumerate}

\subsection{{\bfseries\sffamily SERIOUS} [SR3] The problem of time is unaddressed}
\label{sec:org6231f6e}
The Lagrangian and Hamiltonian use \(\partial N/\partial t\). In a pre-geometric
description, \(t\) must emerge from the substrate. The description currently uses
background time without acknowledging this as a structural gap. This is
particularly acute given the LQG analogy in Doc 4 --- the problem of time is one
of LQG's central unresolved issues.

\subsubsection{Action steps}
\label{sec:org35f8dd9}

\begin{enumerate}
\item{$\square$} Add a brief section to Doc 1 or Doc 4 titled "The Problem of Time"
identifying what \(t\) is in the description. Candidates: a discrete count
of \(\iota\) negotiation steps; a relational time defined by the \(\iota\) cycle
count relative to a reference mode; or a classical background retained as an
approximation. State which is assumed and why.
\end{enumerate}

\subsection{{\bfseries\sffamily SERIOUS} [SR4] The Hamiltonian floor condition is underived}
\label{sec:org23bf7d7}
Several results are conditional on the Hamiltonian being undefined below the
vacuum floor \(\rho_\Lambda\). This is now correctly flagged as conditional in
Docs 3 and 5. But the floor condition is also the only route to addressing the
G and m\_p variation tensions (SEVERE [S1]) and the only route to the Λ hierarchy
resolution. Its derivation is therefore load-bearing for multiple claims.

\subsubsection{Action steps}
\label{sec:org436a0be}

\begin{enumerate}
\item{$\square$} Attempt to derive the floor condition from the substrate Lagrangian: show
that modes below \(\rho_\Lambda\) have zero amplitude in the path integral, or
that the Hamiltonian is ill-defined for configurations below the Hubble-floor
energy density.
\item{$\square$} If derivation is not achievable: consolidate all conditional results into
a single "conditional on floor derivation" section in Doc 1, so the dependency
is visible in one place.
\end{enumerate}

\section{{\bfseries\sffamily MODERATE} Issues}
\label{sec:org81d121e}

\subsection{{\bfseries\sffamily MODERATE} [M1] The LQG minimum area discrepancy has no resolution}
\label{sec:org573e79c}
The 4π factor difference between The Medium's area formula and the LQG formula
is now honestly flagged. But its origin is unknown. A sphere-area normalization
or solid-angle factor is suggested as a candidate; this has not been checked.

\subsubsection{Action steps}
\label{sec:orgebb47fa}

\begin{enumerate}
\item{$\square$} Work out the iota surface area calculation explicitly and determine where
a factor of \(4\pi\) could enter. Either find it or confirm it is absent and
the correspondence is structural only.
\end{enumerate}

\subsection{{\bfseries\sffamily MODERATE} [M2] Metric identification h\_ij (L\^{}\{-1\}) (L(L\^{}\{-1\}))\_ij not justified}
\label{sec:org2dbf99a}
The inverse Laplacian is the graph Green's function, encoding signal propagation,
not a metric in the Riemannian sense. The identification requires showing correct
transformation under graph automorphisms, correct signature, and recovery of flat
space in vacuum.

\subsubsection{Action steps}
\label{sec:org998e2d8}

\begin{enumerate}
\item{$\square$} Either derive the metric properties of \((L^{-1})_{ij}\) explicitly, or
replace the identification with: "The inverse Laplacian is a candidate for
the emergent metric; establishing its Riemannian properties is open work."
\end{enumerate}

\subsection{{\bfseries\sffamily MODERATE} [M3] Barbero-Immirzi derivation remains the single most valuable calculation}
\label{sec:org707f609}
\(\gamma \approx 0.2375\) appears as the residual correction in three separate
gaps: \(G\) (3--6$\backslash$%), proton mass (10$\backslash$%), and \(\hbar\) (0.9$\backslash$%). In LQG,
\(\gamma\) is fixed by black hole entropy counting. If the same entropy counting
performed on the negotiation graph gives \(\gamma \approx 0.2375\), it would
simultaneously close all three gaps and confirm that \(\gamma\) is not a free
parameter. This is the single calculation with the highest payoff across the
entire description.

\subsubsection{Action steps}
\label{sec:orge8ab0ea}

\begin{enumerate}
\item{$\square$} Attempt the Bekenstein-Hawking entropy count on the negotiation graph of
a black hole horizon: count the microstate degeneracy of the negotiation modes
on the horizon graph and require it to equal \(A/4l_\text{Planck}^2\).
\item{$\square$} Even a partial result --- showing the entropy count has the right
structure to fix \(\gamma\) --- would be significant.
\end{enumerate}

\section{Priority Order for Next Working Session}
\label{sec:org83a295b}

In order of what unblocks the most downstream work:

\begin{enumerate}
\item \textbf{FATAL-F2 (text fix):} Remove "from First Principles" from Doc 2 title page.
Takes five minutes. Removes an immediate credibility problem.

\item \textbf{FATAL-F1 (foundational gap):} Connect \(\phi\) to \(N\), or consolidate
around \(N\) as the true primitive. This is the deepest structural gap. Even
a clear statement of what \(N\) is and where it comes from would help.

\item \textbf{MODERATE-M3 (Barbero-Immirzi):} The entropy counting calculation. High
payoff: closes three gaps simultaneously and may constrain the G and m\_p
variation rates if the floor condition falls out of it.

\item \textbf{SEVERE-S1 (G and m\_p variation):} Determine whether a suppression mechanism
exists. If not, formally withdraw those predictions. Leaving them as
"may be ruled out" is better than before but a physicist will want resolution.

\item \textbf{SERIOUS-SR1 (symmetry group):} Complete the group theory calculation for
\(K_3\). This either establishes the SU(3) correspondence or closes the door
on it, which is better than the current uncertainty.

\item \textbf{SERIOUS-SR2 (pion):} Investigate Goldstone / non-complete-graph
accommodations. If neither works, state so explicitly. The pion cannot be
ignored indefinitely.

\item \textbf{SEVERE-S2 (fermionic statistics):} Try the composite antisymmetrization
route. If it works it resolves the most fundamental physical gap in the
description.

\item \textbf{SERIOUS-SR3 and SR4 (time and floor condition):} Add the problem of time
section; attempt the floor derivation. Both are required for the description
to be internally consistent at a foundational level.
\end{enumerate}

\section{What Does Not Need Fixing}
\label{sec:org01c00a3}

\begin{itemize}
\item The internal consistency triangle (\(G\), \(H_0\), \(m_\text{proton}\)) is a
genuine structural observation and should be highlighted, not hedged further.
\item The openbox and plainbox structure across Docs 2--5 is good practice and
should be extended to any new sections rather than removed.
\item The honest labeling of open calculations, failed predictions, and structural
candidates rather than established results is the right register for the
current state of the work. Do not retreat from this.
\item Doc 0 (personal introduction) is in good shape and sets the right tone.
\item The topology hierarchy (\(n=2,3,4,5\) complete graphs) is a genuine structural
result even if incomplete. Keep it.
\end{itemize}
\end{document}
